\documentclass[11pt]{article}
\usepackage{jtsguide}
\setshorttitle{Developer's Guide}

\title{JTS Topology Suite --- Developer's Guide}
\author{Jonathan Aquino (v1.4) \\ rebuilt from the 2003 text}
\date{}

\begin{document}
\jtscover{Developer's Guide}{Version 1.4, rebuilt}{Base text: 17 October 2003 \textperiodcentered\ rebuild: 16 August 2026}

\noindent
The JTS Topology Suite is a Java API for spatial data operations. This
manual keeps the 2003 Developer's Guide chapter plan (getting started,
relate, overlay, buffer, polygonize, line merge, custom coordinates,
tips) and adds a short honest account of the SQL/MM Spatial
(ISO/IEC 13249-3) curve types that are on this tree.

Package names in this tree are \texttt{org.locationtech.jts}. The 2003
copies used \texttt{com.vividsolutions.jts}. LocationTech JTS remains
the upstream project; this fork is not described as if it were already
upstream.

\clearpage
\changectrl{%
  1.4 & 17 Oct 2003 & Jonathan Aquino & Initial draft, created for JTS 1.4. \\
  1.4-rebuild & 16 Aug 2026 & This tree & LaTeX rebuild of the 2003 guide.
    Curve, I/O, overlay, and Hausdorff claims match the code on this
    tree. ACME licence text is not the project licence. \\
}

\tableofcontents
\clearpage

\section{Overview}
The JTS Topology Suite is a Java API that implements a core set of
spatial data operations using an explicit precision model and robust
geometric algorithms. It provides a complete model for specifying 2-D
linear Geometry. Many common operations in computational geometry and
spatial data processing are exposed in a clear, consistent and
integrated API. JTS is intended to be used in the development of
applications that support the validation, cleaning, integration and
querying of spatial datasets.

This document is intended for developers who would like to use JTS to
accomplish their spatial data processing requirements. It describes
common uses of the JTS API and gives code examples.

The 2003 note ``this document is under construction'' is left as
history. The linear API walkthrough is still the right starting point.
An opt-in curve module sits beside that model; it does not replace it.

\subsection{Other resources}
\begin{itemize}
  \item OpenGIS Simple Features Specification For SQL Revision 1.1
        (SFS). The reference specification for the linear spatial data
        model and the spatial predicates and functions implemented by
        JTS.
  \item SQL/MM Spatial (ISO/IEC 13249-3) --- CircularString,
        CompoundCurve, CurvePolygon, and WKB codes 8--12 used by the
        opt-in curve module. No 13249-3 DOI is cited.
  \item \emph{JTS Technical Specifications}. The design specification
        for the classes, methods and algorithms.
  \item JTS JavaDoc for packages under \texttt{org.locationtech.jts}.
\end{itemize}

\section{Getting started}
The most common JTS tasks involve creating and using Geometry objects.
The easiest way to create a linear Geometry by hand is to use a
\texttt{WKTReader} to generate one from a Well-Known Text (WKT)
string. For example:

\begin{lstlisting}
Geometry g1 = new WKTReader().read(
    "LINESTRING (0 0, 10 10, 20 20)");
\end{lstlisting}

A precise specification for WKT is given in the JTS Technical
Specifications. Many examples of WKT may be found in the files in the
test directory.

In a real program, it is easier to use a \texttt{GeometryFactory},
because you do not need to build up a WKT string; rather, you work
with the objects directly:

\begin{lstlisting}
Coordinate[] coordinates = new Coordinate[] {
    new Coordinate(0, 0), new Coordinate(10, 10),
    new Coordinate(20, 20) };
Geometry g1 = new GeometryFactory().createLineString(coordinates);
\end{lstlisting}

Once you have a Geometry you can compute intersection, area, envelope,
centroid, and buffer. See the JavaDoc for
\texttt{org.locationtech.jts.geom.Geometry}.

\begin{amendment}
On this tree the concrete SQL/MM Spatial (ISO/IEC 13249-3) curve
types live in the
\texttt{jts-curve} module
(\texttt{org.locationtech.jts.geom.curve}):
\texttt{CircularString}, \texttt{CompoundCurve},
\texttt{CurvePolygon}, and the collections \texttt{MultiCurve} and
\texttt{MultiSurface}. Construction is opt-in: the default
\texttt{GeometryFactory} throws on \texttt{createCircularString} /
\texttt{createCompoundCurve} / \texttt{createCurvePolygon} /
\texttt{createMultiCurve} / \texttt{createMultiSurface}. Use
\texttt{CurveGeometryFactory}. Constructors do not linearise.

A core \texttt{WKTReader} still throws \texttt{ParseException} on
\texttt{CIRCULARSTRING}. Opt in:
\end{amendment}

\begin{lstlisting}
CurveGeometryFactory gf = new CurveGeometryFactory();
CurveWKTReader rdr = new CurveWKTReader(gf);
Geometry arc = rdr.read("CIRCULARSTRING (0 0, 5 5, 10 0)");
// arc.getGeometryType() == "CircularString"
\end{lstlisting}

\figshot{ug1_dg1_draw_cs.png}{DG-1. Draw CircularString.}

\section{Computing spatial relationships}
An important application of JTS is computing the spatial relationships
between Geometries. JTS follows the Dimensionally Extended
9-Intersection Matrix (DE-9IM) model specified by the OGC. To compute
the DE-9IM for two Geometries, use the \texttt{relate} method:

\begin{lstlisting}
Geometry a = /* ... */;
Geometry b = /* ... */;
IntersectionMatrix m = a.relate(b);
\end{lstlisting}

Most relationships of interest can be specified as a pattern which
matches a set of intersection matrices. JTS also provides a set of
boolean predicates which compute common spatial relationships
directly:

\begin{center}
\begin{tabular}{>{\raggedright\arraybackslash}p{2.6cm}p{10.6cm}}
  \toprule
  \textbf{Method} & \textbf{Meaning} \\
  \midrule
  Equals     & The Geometries are topologically equal. \\
  Disjoint   & The Geometries have no point in common. \\
  Intersects & The Geometries have at least one point in common
               (the inverse of Disjoint). \\
  Touches    & The Geometries have at least one boundary point in
               common, but no interior points. \\
  Crosses    & The Geometries share some but not all interior points,
               and the dimension of the intersection is less than that
               of at least one of the Geometries. \\
  Within     & Geometry A lies in the interior of Geometry B. \\
  Contains   & Geometry B lies in the interior of Geometry A
               (the inverse of Within). \\
  Overlaps   & The Geometries share some but not all points in common,
               and the intersection has the same dimension as the
               Geometries themselves. \\
  \bottomrule
\end{tabular}
\end{center}

In some cases the precise definition of the predicates is subtle.
Refer to the JTS Technical Specifications. On curve operands, public
\texttt{relate} / predicates that have not taken a certified-disc path
still see control-point chords unless the caller linearises explicitly
via \texttt{Linearizable.toLinear(tolerance)}.

\section{Computing overlay operations}
The previous section discussed functions that return true or false,
like Intersects and Contains. The JTS overlay operations, some of
which the 2003 guide illustrated as a six-panel cartoon, are:

\begin{center}
\begin{tabular}{>{\raggedright\arraybackslash}p{3.2cm}p{10.0cm}}
  \toprule
  \textbf{Method} & \textbf{Meaning} \\
  \midrule
  Buffer         & The Polygon or MultiPolygon which contains all
                   points within a specified distance of the Geometry.
                   See Section~\ref{sec:buffers}. \\
  ConvexHull     & The smallest convex Polygon that contains all the
                   points in the Geometry. \\
  Intersection   & The set of all points which lie in both A and B. \\
  Union          & The set of all points which lie in A or B. \\
  Difference     & The set of all points which lie in A but not in B. \\
  SymDifference  & The set of all points which lie in either A or B
                   but not both. \\
  \bottomrule
\end{tabular}
\end{center}

As with the spatial relationships, these operations have precise
definitions given in the JTS Technical Specifications.

On linear Geometry the production overlay is OverlayNG
(\texttt{org.locationtech.jts.operation.overlayng}). The 2003
noding-graph overlay is historical.

\begin{amendment}
The curve overlay type is \texttt{OverlayNGCurve}. It lives in
\texttt{jts-curve}
(\texttt{org.locationtech.jts.operation.overlayng.curve}). There is
no public noder. Closed-form answers are limited to certified discs.
Anything else densifies and delegates to core OverlayNG.
\texttt{Geometry.intersection} / \texttt{union} / \texttt{difference}
/ \texttt{symDifference} on a curve type are not a general circular
overlay.

Closed-form constructions on this tree are certified discs, the
single-arc hull, and the compound-curve hull of circular-plus-straight
members (including a stadium). Disc buffer of a full circular disc is
a disc of radius $r+d$ (or empty if $r+d \le 0$). An open arc stays
on \texttt{BufferOp} chords. Anything that is not a certified closed
form uses \texttt{toLinear} and the linear algorithm --- a fallback,
not a fail, and not a closed form.
\end{amendment}

\begin{lstlisting}
import org.locationtech.jts.operation.overlayng.curve.OverlayNGCurve;
Geometry cap = OverlayNGCurve.intersection(a, b);
Geometry cup = OverlayNGCurve.union(a, b);
Geometry sub = OverlayNGCurve.difference(a, b);
Geometry xor = OverlayNGCurve.symDifference(a, b);
\end{lstlisting}

A certified disc of radius 5 has convex-hull area $25\pi$. That hull
is the disc itself, not a densified polygon.

\figshot{ug7_dg3_ts2_disc_hull.png}{DG-3. Disc hull area $25\pi$. Circular-plus-straight laser hull shipped.}

\section{Computing buffers}
\label{sec:buffers}
In GIS, buffering is an operation used to compute the area containing
all points within a given distance of a Geometry. In mathematical
terms, this is computing the Minkowski sum of the Geometry with a
disc of radius equal to the buffer distance. Finding positive and
negative buffers is sometimes referred to as erosion and dilation. In
CAD/CAM buffer curves are called offset curves.

You can use JTS to compute the buffer of a Geometry using the
\texttt{Geometry.buffer} method or the \texttt{BufferOp} class. The
input Geometry may be of any type (including arbitrary
GeometryCollections). The result of a buffer operation is always an
area type (Polygon or MultiPolygon). The result may be empty (for
example, a negative buffer of a LineString).

You can compute buffers with both positive and negative buffer
distances. Buffers with a positive buffer distance always contain the
input Geometry. Buffers with a negative buffer distance are always
contained within the input Geometry. A negative buffer of a
LineString or a Point results in an empty Geometry.

Buffer distances of 0 are also supported. You can use this to perform
an efficient union of multiple polygons.

\subsection{Basic buffering}
To compute a buffer for a given distance, call \texttt{buffer()} on
the Geometry:

\begin{lstlisting}
Geometry g = /* ... */;
Geometry buffer = g.buffer(100.0);
\end{lstlisting}

\subsection{End cap styles}
Buffer polygons can be computed with different line end-cap styles.
The end-cap style determines how the linework for the buffer polygon
is constructed at the ends of linestrings.

\begin{center}
\begin{tabular}{>{\raggedright\arraybackslash}p{3.6cm}p{9.6cm}}
  \toprule
  \textbf{Style} & \textbf{Description} \\
  \midrule
  CAP\_ROUND          & The usual round end caps. \\
  CAP\_FLAT / CAP\_BUTT & End caps truncated flat at the line ends
                         (the 2003 name was CAP\_BUTT). \\
  CAP\_SQUARE         & End caps squared off at the buffer distance
                         beyond the line ends. \\
  \bottomrule
\end{tabular}
\end{center}

To specify the buffer end-cap style, the \texttt{BufferOp} class in
\texttt{org.locationtech.jts.operation.buffer} is used directly:

\begin{lstlisting}
Geometry g = /* ... */;
BufferOp bufOp = new BufferOp(g);
bufOp.setEndCapStyle(BufferParameters.CAP_FLAT);
Geometry buffer = bufOp.getResultGeometry(distance);
\end{lstlisting}

\subsection{Specifying the approximation quantization}
Since the exact buffer outline of a linear Geometry usually contains
circular sections, the buffer must be approximated by the linear
Geometry supported by JTS. The degree of approximation may be
controlled by specifying the number of quadrant segments used to
approximate a quarter-circle. Specifying a larger number of segments
results in a better approximation to the actual area, but also
results in a larger number of line segments in the computed polygon.

The default number of segments is 8. This gives less than a 2\%
maximum error in the distance of the computed curve approximation to
the actual buffer curve. This error can be reduced to less than 1\%
by using a value of 12.

\begin{lstlisting}
Geometry buffer = g.buffer(100.0, 16);
\end{lstlisting}

\begin{amendment}
A full circular disc (certified \texttt{CurvePolygon}) has a
closed-form buffer: another disc of radius $r+d$, or empty if
$r+d \le 0$. An open \texttt{CircularString} is not a disc. Its
buffer is \texttt{BufferOp} on chords. That is not a closed form.
\end{amendment}

\section{Polygonization}
Polygonization is the process of forming polygons from linework which
encloses areas. Linework to be formed into polygons must be fully
noded --- that is, linestrings must not cross and must touch only at
endpoints.

JTS provides the \texttt{Polygonizer} class to perform
polygonization. The Polygonizer takes a set of fully noded
LineStrings and forms all the polygons which are enclosed by the
lines. Polygonization errors such as dangling lines or cut lines can
be identified and reported.

\begin{lstlisting}
Polygonizer polygonizer = new Polygonizer();
polygonizer.add(lines);
Collection polys   = polygonizer.getPolygons();
Collection dangles = polygonizer.getDangles();
Collection cuts    = polygonizer.getCutEdges();
\end{lstlisting}

If the set of lines is not correctly noded the Polygonizer will still
operate on them, but the resulting polygonal Geometries will not be
valid. The MultiLineString union technique can be used to node a set
of LineStrings (see Section~\ref{sec:noding}).

\section{Merging a set of LineStrings}
Sometimes a spatial operation such as union will produce chains of
small LineStrings. The JTS \texttt{LineMerger} is a simple utility to
sew these small LineStrings together.

The LineMerger assumes that the input LineStrings are noded (they do
not cross; only their endpoints can touch). The output LineStrings
are also noded. If LineStrings to be merged do not have the same
direction, the direction of the resulting LineString will be that of
the majority.

\begin{lstlisting}
LineMerger lineMerger = new LineMerger();
lineMerger.add(lineStrings);
Collection merged = lineMerger.getMergedLineStrings();
\end{lstlisting}

\section{Using custom CoordinateSequences}
By default JTS uses arrays of Coordinates to represent the points and
lines of Geometries. There are some cases in which you might want
Geometries to store their points using some other implementation. For
example, to save memory you may want to use a more compact sequence
implementation, such as an array of $x$'s and an array of $y$'s.
Another possibility is to use a custom coordinate class to store
extra information on each coordinate, such as measures for linear
referencing.

You can do this by implementing the \texttt{CoordinateSequence} and
\texttt{CoordinateSequenceFactory} interfaces. You would then create
a \texttt{GeometryFactory} parameterized by your
\texttt{CoordinateSequenceFactory}, and use this GeometryFactory to
create new Geometries. All of these new Geometries will use your
CoordinateSequence implementation.

If your CoordinateSequence is not based on an array of the standard
JTS Coordinates (or a subclass of Coordinate), it may incur a small
performance penalty due to marshalling. Curve types use the same
factory hook; \texttt{CurveGeometryFactory} can be constructed with a
custom \texttt{CoordinateSequenceFactory}.

\section{Tips and techniques}

\subsection{Noding a set of LineStrings}
\label{sec:noding}
Many spatial operations assume that their input data is noded,
meaning that LineStrings never cross. For example, the JTS
Polygonizer and the JTS LineMerger assume that their input is noded.

The noding process splits LineStrings that cross into smaller
LineStrings that meet at a point, or node. A simple trick for noding
a group of LineStrings is to union them together; the unioning
process will node the LineStrings.

\begin{lstlisting}
Collection lineStrings = /* ... */;
Geometry noded = (LineString) lineStrings.get(0);
for (int i = 1; i < lineStrings.size(); i++) {
    noded = noded.union((LineString) lineStrings.get(i));
}
\end{lstlisting}

Unary union / OverlayNG is the current production path for many
polygons. The 2003 ``buffer of a GeometryCollection at distance 0''
trick still works but is no longer the first recommendation.

\subsection{Unioning many polygons efficiently}
Calling \texttt{Polygon.union} repeatedly is one way to union several
Polygons together. Prefer \texttt{UnaryUnionOp} / OverlayNG union on
a collection. The 2003 \texttt{GeometryCollection.buffer(0)} trick
remains documented as history.

\section{Curve types on this tree}
\label{sec:curves}
On this tree the concrete SQL/MM Spatial (ISO/IEC 13249-3) curve
types live in the
\texttt{jts-curve} module
(\texttt{org.locationtech.jts.geom.curve}):

\begin{center}
\begin{tabular}{lll}
  \toprule
  \textbf{WKT keyword} & \textbf{Java class} & \textbf{Extends} \\
  \midrule
  CIRCULARSTRING & \texttt{...geom.curve.CircularString} & LineString \\
  COMPOUNDCURVE  & \texttt{...geom.curve.CompoundCurve}  & LineString \\
  CURVEPOLYGON   & \texttt{...geom.curve.CurvePolygon}   & Polygon \\
  MULTICURVE     & \texttt{...geom.curve.MultiCurve}     & MultiLineString \\
  MULTISURFACE   & \texttt{...geom.curve.MultiSurface}   & MultiPolygon \\
  \bottomrule
\end{tabular}
\end{center}

Maven artifact: \texttt{org.locationtech.jts:jts-curve}. Callers
instantiate \texttt{CurveGeometryFactory} / \texttt{CurveWKTReader} /
\texttt{CurveWKBWriter} explicitly. The module does not auto-register
via ServiceLoader.

\texttt{Linearizable.toLinear(tolerance)} is the explicit densify.
Spatial methods that have not taken a certified closed form still
fall through to the parent linear type or to a linearised copy. That
is not hidden linearise-on-construct.

Closed-form constructions on this tree:
\begin{itemize}
  \item Certified disc: area, envelope, disc-buffer, disc hull
        (radius-5 disc hull area $25\pi$).
  \item Single-arc convex hull: the slice bounded by the arc and its
        chord, kept as a \texttt{CurvePolygon}.
  \item Compound-curve hull of circular-plus-straight members
        (H-CC), including a stadium. The exposed arc stays a
        \texttt{CircularString}; the result is a
        \texttt{CurvePolygon}, not a densified polygon. A clothoid
        hull stays the linear fallback
        (\texttt{CurveExact} null $\rightarrow$
        \texttt{linearise().convexHull()}), not a laser.
  \item \texttt{MaximumInscribedCircle} on a certified disc and on a
        stadium; \texttt{LargestEmptyCircle} on legal curve
        obstacles. These are constructions, not a general circular
        noder.
\end{itemize}

The curve overlay type is \texttt{OverlayNGCurve}. There is no public
noder. Closed-form overlay answers are limited to certified discs.
TestBuilder can draw \texttt{logoLines} as curves.

\section{WKT and WKB on this tree}
Core \texttt{WKTReader} still rejects \texttt{CIRCULARSTRING} and the
other SQL/MM keywords. Read them with \texttt{CurveWKTReader}. Write
structured members with \texttt{CurveWKTWriter}. Core
\texttt{WKTWriter} already emits the subclass keyword via
\texttt{getGeometryType()} for a flat \texttt{CircularString};
composite types need the curve writer to keep member tags.

\texttt{CLOTHOID} is recognised only by \texttt{CurveWKTReader}, and
only as a non-leading \texttt{COMPOUNDCURVE} member. That path is
extras, not a closed-form construction.

WKB type codes 8--12 are first-class in core \texttt{WKBReader} and
\texttt{WKBConstants} (\texttt{CircularString}=8,
\texttt{CompoundCurve}=9, \texttt{CurvePolygon}=10,
\texttt{MultiCurve}=11, \texttt{MultiSurface}=12). Construction still
requires a curve-capable factory; otherwise the reader throws.
\texttt{CurveWKBWriter} emits codes 8--12. A bare
\texttt{WKBWriter} still walks the parent \texttt{instanceof} ladder
and writes a \texttt{CircularString} as LineString (code 2) and a
\texttt{CurvePolygon} as Polygon (code 3). That is what a direct
\texttt{WKBWriter.write} call does on this tree.

Default and export WKB paths (JtsOp, TestBuilder XML export) route
through \texttt{CurveWKBWriter} so those writers do not flatten
codes 8--12.

\begin{lstlisting}
// Read codes 8-12 as curve types:
Geometry g = new WKBReader(new CurveGeometryFactory()).read(bytes);
Geometry g2 = new CurveWKBReader().read(bytes);

// Write codes 8-12:
byte[] wkb = new CurveWKBWriter().write(g);

// Bare writer still emits 2 / 3 for the parent types:
byte[] flat = new WKBWriter().write(g); // CircularString -> type 2
\end{lstlisting}

\fighole{DG-2 WKB type 8 (CurveWKBWriter)}

\section{Hausdorff on this tree}
The 2003 Developer's Guide did not document Hausdorff. It is recorded
here so a reader is not left with the public-API chord behaviour as
an unspoken surprise.

Public \texttt{DiscreteHausdorffDistance} on this tree (commit
\texttt{0ca71b}) has a closed form for two pairs only:
\begin{enumerate}
  \item A single-arc \texttt{CircularString} toward a single-segment
        \texttt{LineString}, apex $\sqrt{949}/6 - 7/6 = 3.967640600249787$
        for the witness
        \texttt{CIRCULARSTRING (0 0, 2 3, 10 0)} vs
        \texttt{LINESTRING (0 0, 10 0)}.
  \item Two circular discs (10.0). A
        \texttt{MultiSurface} unwrap of one disc is the same disc
        pair, not a third.
\end{enumerate}
The exact path skips densify; densify is not the laser. For every
other pair the public API still sees vertices / chords. Tests keep
\texttt{fail()} on the general case. Fr\'echet remains open.

\fighole{DG-4 D-HF discs 10.0}

\section{Licence}
JTS is dual-licensed under the Eclipse Public License 2.0 and the
Eclipse Distribution License 1.0. See \texttt{LICENSE\_EPLv2.txt},
\texttt{LICENSE\_EDLv1.txt}, and \texttt{LICENSES.md}. The 2005
TestBuilder guide's ACME licence appendix is historical and is not
the project licence.

\vspace{1.2em}
\noindent\textit{End of Developer's Guide. Remaining named holes
(WKB, D-HF) are empty 4:3 slots (target 1600$\times$1200, full
TestBuilder window). 2003 Amyuni conversion cartoons are omitted;
the prose they illustrated is kept.}

\end{document}
